\documentclass[12pt]{amsart}
\addtolength{\topmargin}{-0.6in} % usually -0.25in
\addtolength{\textheight}{1.1in} % usually 1.25in
\addtolength{\oddsidemargin}{-0.7in}
\addtolength{\evensidemargin}{-0.7in}
\addtolength{\textwidth}{1.5in} %\setlength{\parindent}{0pt}

\newcommand{\normalspacing}{\renewcommand{\baselinestretch}{1.04}\tiny\normalsize}

\normalspacing

% macros
\usepackage{amssymb,xspace}
\usepackage{tikz}
\usepackage[pdftex,colorlinks=true]{hyperref}


\newtheorem*{thm}{Theorem}
\newtheorem*{lem}{Lemma}

\newcommand{\mtt}{\texttt}
\newcommand{\mtl}[1]{{\texttt{>>#1}}}
\usepackage{alltt}
\usepackage{fancyvrb}

\newcommand{\bu}{\mathbf{u}}
\newcommand{\bv}{\mathbf{v}}

\newcommand{\CC}{{\mathbb{C}}}
\newcommand{\RR}{{\mathbb{R}}}
\newcommand{\ZZ}{{\mathbb{Z}}}
\newcommand{\ZZn}{{\mathbb{Z}}_n}
\newcommand{\NN}{{\mathbb{N}}}

\newcommand{\eps}{\epsilon}
\newcommand{\grad}{\nabla}
\newcommand{\lam}{\lambda}
\newcommand{\ip}[2]{\mathrm{\left<#1,#2\right>}}
\newcommand{\erf}{\operatorname{erf}}

\renewcommand{\Re}{\operatorname{Re}}
\renewcommand{\Im}{\operatorname{Im}}
\newcommand{\Arg}{\operatorname{Arg}}

\newcommand{\Span}{\operatorname{span}}
\newcommand{\rank}{\operatorname{rank}}
\newcommand{\range}{\operatorname{range}}
\newcommand{\trace}{\operatorname{tr}}
\newcommand{\Null}{\operatorname{null}}

\newcommand{\Matlab}{\textsc{Matlab}\xspace}
\newcommand{\Octave}{\textsc{Octave}\xspace}
\newcommand{\pylab}{\textsc{pylab}\xspace}
\newcommand{\longMOP}{\textsc{Matlab}\big|\textsc{Octave}\big|\textsc{pylab}\xspace}
\newcommand{\MOP}{\textsc{M}\big|\textsc{O}\big|\textsc{p}\xspace}

\newcommand{\prob}[1]{\bigskip\noindent\large\textbf{#1.} \normalsize}
\newcommand{\bookprob}[1]{\bigskip\noindent\large\textbf{Exercise #1.} \normalsize}
\newcommand{\probpart}[1]{\smallskip\noindent\textbf{(#1)}\quad }
\newcommand{\aprobpart}[1]{\textbf{(#1)}\quad }


\newcommand*\circled[1]{\tikz[baseline=(char.base)]{
            \node[shape=circle,draw,inner sep=2pt] (char) {#1};}}


\begin{document}
\scriptsize \noindent Math 617 Functional Analysis (Bueler) \hfill assigned: 22 April 2026
\thispagestyle{empty}

\bigskip
\Large\textbf{\centerline{Final Exam: Prove 3 Nontrivial Theorems}}

\medskip
\large\textbf{\centerline{Friday, 1 May, 10:15am--12:15pm, Chapman 107}}

\normalsize
\bigskip
For the in-class Final Exam I would like you to \textbf{prove 3 nontrivial theorems} which we have seen during the semester.  A list of 9 possibilities appears below, and these are the only possibilities.  You will have this document in your hand during the Exam, but you must remember what you want to say!  \textbf{You may NOT bring notes, the textbook, or electronics of any kind to the Exam.}  Just bring a writing implement.

It is essential that you \textbf{draft and practice both theorem statements and the proofs}.  Make sure to read relevant sections of the textbook or slides.  Get feedback on your drafts from other students or friends (or me).  Decide in advance how you will remember enough detail so as to recreate the theorem statement and the proof during the Exam itself.

\textbf{This is a writing assignment.}  I will grade for correctness, completeness, and readability, in that order of priority.

\medskip
\noindent \hrulefill

\noindent \textbf{\large Instructions:}

\smallskip
Write each theorem on a separate sheet.  Put your name in the upper-right corner.  Identify each theorem by ``\textbf{T1}'', etc.~in the upper-left corner.  The statement of the theorem and the full proof should be at most 1.5 pages.  Clearly and precisely state the theorem, making sure that both the hypotheses and conclusion are correct.  Then write the proof.  Full credit will only be given for a well-written, appropriately-detailed, and correct proof.  I recommend starting your proof with a sketch of its strategy, because substantial partial credit can then be given even if certain details are incomplete or incorrect.

\medskip
\noindent \hrulefill

\medskip
\renewcommand{\labelenumi}{\textbf{T\arabic{enumi}.}}
\setlength{\leftmargini}{26pt}
\begin{enumerate}
\setlength{\itemsep}{4pt}
\item xx
\end{enumerate}

\end{document}

